% Generated draft: requires mathematical and source review.
\documentclass[11pt]{article}
\usepackage[margin=25mm]{geometry}
\usepackage{fontspec}
\setmainfont{Latin Modern Roman}
\usepackage{amsmath,amssymb,mathtools,unicode-math}
\setmathfont{Latin Modern Math}
\usepackage{enumitem,needspace,xurl,hyperref}
\setlength{\parindent}{0pt}
\setlength{\parskip}{5pt}
\setlength{\emergencystretch}{3em}
\newcommand{\N}{\mathbb{N}}
\newcommand{\Z}{\mathbb{Z}}
\newcommand{\Q}{\mathbb{Q}}
\newcommand{\R}{\mathbb{R}}
\newcommand{\C}{\mathbb{C}}
\newcommand{\F}{\mathbb{F}}
\newcommand{\A}{\mathbb{A}}
\newcommand{\PP}{\mathbb{P}}
\newcommand{\E}{\mathbb{E}}
\newcommand{\eps}{\varepsilon}
\newcommand{\dd}{\,\mathrm{d}}
\begin{document}
\section*{UnsolvedMath ID: 30001668}
\textbf{Source problem code:} catalogue entry
\subsection{Definitions and mathematical statement}
For a finite graph \(G=(V,E)\) of order \(n=|V|\geq 1\), a \emph{prime labeling} is a bijection
\[
f:V\longrightarrow\{1,\ldots,n\}
\]
such that
\[
\gcd\bigl(f(u),f(v)\bigr)=1
\qquad\text{for every }\{u,v\}\in E.
\]
The graph is called \emph{prime} if it admits such a labeling; this terminology does not require the individual vertex labels to be prime numbers.

Using the standard convention that a tree is a nonempty finite simple undirected connected acyclic graph, the supplied question is:
\[
\begin{gathered}
\text{For every integer }n\geq 1\text{ and every tree }T=(V,E)\text{ with }|V|=n,\\
\text{does there exist a bijection }f:V\longrightarrow\{1,\ldots,n\}\\
\text{such that }\gcd\bigl(f(u),f(v)\bigr)=1
\text{ for every }\{u,v\}\in E?
\end{gathered}
\]
For \(n=1\), the adjacency condition is vacuous. The source does not explicitly specify graph conventions or whether an empty tree is permitted; the convention used above is therefore stated explicitly.

The supplied record labels the problem \texttt{partially\_solved}, but its accompanying literature assessment concerns the graceful tree conjecture and graceful labelings. That assessment does not establish the status of the prime-labeling question stated here. The intended connection, if any, remains unresolved.
\subsection{Short English statement}
Can the vertices of every finite tree with \(n\) vertices be assigned the integers \(1,\ldots,n\), each exactly once, so that the labels at the ends of every edge are coprime? The supplied literature assessment discusses graceful labelings, so its relevance to this question requires review.
\subsection{Sources}
\begin{itemize}
\item Supplied local UnsolvedMath registry record, permanent ID \(30001668\), title \emph{Prime Labelings of Trees}, exact source problem code \texttt{catalogue entry}, in the Oberwolfach Reports Open Problems collection. The record gives proposed year 2011 and supplies no named proposer.
\item Supplied report citation: \emph{Combinatorics} (2011), Oberwolfach Reports, pp.~5--83, DOI \texttt{10.4171/owr/2011/01}. Source URL and literature-source label ``Oberwolfach problem statement'': \url{https://doi.org/10.4171/owr/2011/01}.
\item Supplied literature source labeled ``Gallian survey'': \url{https://www.combinatorics.org/files/Surveys/ds6/ds6v8-2003.pdf}.
\item Supplied literature source labeled ``current conjecture status'': \url{https://mlelarge.github.io/graph-conjectures/op/graceful_tree_conjecture/}.
\item Literature assessment checked at \texttt{2026-08-21}: this is supplied dated curation derived from the OWR curation, not an independently verified assessment. It labels the record partially solved and states that the graceful tree conjecture remains open in general while many special tree classes, including ``several listed here,'' are known to be graceful. No such list is present in the supplied record, and the assessment concerns a different labeling property from the stated prime-labeling question. The source itself requests independent verification of time-sensitive status and cited claims.
\item The supplied statement status is \texttt{corrected\_verified}; its verification note describes only enclosing TeX notation in mathematical delimiters without changing mathematical content. This supplied label is not independent mathematical or source verification.
\end{itemize}
\end{document}
